\pdfbookmark[1]{How this was found}{Provenance}
\chapter*{How this was found}

This section records the order in which things happened, because the order bears on how the results
should be read.

\section*{The author's position}

I have no classical training in mathematics, linguistics, crystallography, cryptography or
information theory. I am self-taught across all of them, arrived at each one sideways, and did not
come to any of them through a curriculum, a supervisor or a department. I also have difficulty
communicating, and a tool checks the prose in these books instead of my trusting it.

That is stated first because it explains the shape of what follows. Somebody trained in a field
reads its literature before building, and meets a precedent before writing a line. I built first and
read afterward. Every precedent named in these books was found that way, and finding one late is not
the same as finding one early.

\section*{Why the refuted entries are kept}

Honesty is the stated reason and it is not the working one.

I do not limit myself to conventional thought, or to conventional mathematics, and I do not accept a
result because it arrived looking the way a result is supposed to look. That is how several things
here were found. It is also how somebody working this way gets a long way down a wrong road without
anybody stopping them, because there is no field standing behind me to say the road is wrong.

So the workbook is the thing that stops me. Writing down what a measurement killed, beside the claim
it killed, is how I keep scrutinizing what I already think I know. An entry that says what was
believed, what was measured, and which of the two lost is a check I can run on myself later, and it
is the only check available to somebody with no supervisor and no committee.

So an entry here is not a confession and should not be read as one. Refuted entries are
working notes. The document would be less useful without them and so would I.

\section*{The timeline is in the repository and it is signed}

The first commit of this repository is dated 2026-09-04 and is signed with my GPG key. Its subject is
human language signature detection research. Seven days separate it from the date on this page.

Most of what is cited in these books I had not read before that commit. That is checkable and not
asserted: the history records when each result appeared, when each search for a precedent was
run, and when the citation was entered. A reader who doubts the order can read it off the commits
instead of taking this section's word for it.

The seven days are also why the precedents arrive where they do. There was no interval in which to
read a field first. Each result was built, and then somebody went looking for who had already had
it.

\section*{The order the precedents were found in}

The construction was reached from one general question, about how far an arrangement sits from the
most disordered arrangement of its own parts. No crystal, no cipher and no corpus was in mind. Each
domain here is that question with a different answer to what counts as a part.

Arriving from there at something a field already uses is the good outcome. It is evidence the
general question points at something real, and arriving somewhere nobody had ever needed to go would
have been the weaker result. It gives up priority. I claim none here.

Each precedent was found by asking, once a result existed, how such a thing could have gone
unnoticed. That question is what produced the search, and the search is what produced the citation.

\begin{itemize}
  \item \textbf{The cross-language ordering result.} Montemurro and Zanette,
    \_Universal Entropy of Word Ordering Across Linguistic Families\_, PLoS ONE 6(5):e19875, 2011.
    Two hundred times the sample here, reaching Old Egyptian and Sumerian, and fourteen years
    earlier. The conclusion is theirs. This work reached it independently and later.
  \item \textbf{The null permutation.} Fisher, 1935, and standard in statistics ever since.
  \item \textbf{The maximum entropy reference.} Jaynes, 1957, and standard in physics.
  \item \textbf{Filtering with a necessary condition.} Standard in string matching, with Navarro
    reporting it dominating a large range of parameters. A necessary condition used to discard
    candidates, with a mandatory verification pass, is the filtering paradigm and it has a
    literature. No claim of novelty for the construction survives that reading.
  \item \textbf{The injectivity of the $\Sigma$ functions.} Handschuh and Gilbert, by a polynomial
    argument, reproduced here for all four functions instead of two.
\end{itemize}

\section*{What I have not read}

The crystallographic literature has not been worked through. An autocorrelation of density in a
crystal is very unlikely to be new, a century of the field has worked on exactly that object, and
the reasonable prior is that the construction has a name, a date and a literature behind it. Nobody
here has gone and found out. Until somebody does, nothing in that book is a claim to a new method
and nothing in it is an attribution either.

\section*{What now checks this}

Novelty is a claim about other people's work. Nothing in this repository can check one.

\texttt{maint/\allowbreak{}citations/\allowbreak{}prior\_\allowbreak{}works.\allowbreak{}py} reports
every novelty claim in the books with no citation attached and no line saying I have not read the
literature. I decide each one. No pattern separates my priority claim from my description of
somebody else's work.

\texttt{maint/\allowbreak{}citations/\allowbreak{}citations.\allowbreak{}py} is the other half. A
surname in a comment is not a citation. I register a source once I have it open and can enter an
author, a year, a title and an identifier. An empty row means I have not looked it up yet.
\cleardoublepage
